% ================================================================
\section{The Gram Matrix as a Quantum Correlator}
% ================================================================

\subsection{The Two-Point Function}

The Gram matrix $G \in \RR^{(N-1) \times (N-1)}$ with entries
\[
  G(j,k) = \int_0^1 \fract{1/(jx)}\,\fract{1/(kx)}\,dx
  = \braket{h_j}{h_k}
\]
is the \emph{two-point correlation function} (propagator) of the
prime number quantum field. In condensed matter physics, this is the
Green's function $G(j,k) = \expect{a_j^\dagger a_k}$ measuring the
amplitude for a particle created at site $k$ to propagate to site $j$.

\begin{observation}[Diagonal = Self-Energy (Proved)]
The diagonal entries $G(k,k) = \int_0^1 \fract{1/(kx)}^2\,dx$ are
the \emph{self-energy} of the $k$-th mode. The Cathedral proves
this identity as a \emph{theorem} (April 20, 2026) via a three-step
argument that mirrors the self-energy computation in lattice QFT:
\begin{enumerate}
  \item \textbf{Lattice discretization.} Partition $(0,1)$ into intervals
    where $\lfloor 1/(kx) \rfloor = m$ is constant. On each interval,
    the integrand $\fract{1/(kx)}^2 = (1/(kx) - m)^2$ is a rational
    polynomial whose antiderivative is computed exactly via FTC.
    This is the number-theoretic analogue of a lattice field theory
    with sites indexed by $m$.
  \item \textbf{Thermodynamic limit.} The sum over lattice sites is
    identified with Stirling's partial sum $P(K)$, using the identity
    $P(K) = 2\log K! - 2K\log K + 2K - 1 - H_K$.
    Mathlib's Stirling approximation ($\sqrt{2\pi K}(K/e)^K \to K!$)
    provides the $\ln(2\pi)$ contribution---the thermodynamic
    entropy of the lattice.
  \item \textbf{Mass renormalization.} The Euler--Mascheroni constant
    $\gamma = \lim_{K\to\infty}(H_K - \ln K)$ appears as the ``bare
    mass'' subtraction. A squeeze theorem argument takes the lattice
    spacing to zero: $\int_0^{1/K} \fract{1/u}^2\,du \to 0$.
\end{enumerate}
The result, $G(k,k) = k^{-2}(\ln 2\pi - \gamma - 1)$, is proved from
Mathlib's Stirling and harmonic number infrastructure with \emph{zero
axioms}. The self-energy saturates at $1 - 1/k + O(\ln k / k)$ for
large $k$, consistent with asymptotic freedom.
\end{observation}

\begin{observation}[Off-Diagonal = Interaction (Structurally Proved)]
The off-diagonal entries $G(j,k)$ for $j \neq k$ encode the
\emph{interaction} between modes. The Vasyunin formula expresses these
as cotangent sums---number-theoretic analogues of scattering phases.
The interaction depends on $\gcd(j,k)$, reflecting the
\emph{arithmetic structure} of the coupling: modes interact strongly
precisely when they share prime factors.

The identity $G(j,k) = \text{vasyuninGramFormula}(j,k)$ is now
\textbf{structurally proved} (April~25, 2026) via a
\emph{uniqueness-of-limits} argument: the partial integral
$\int_{1/(jM)}^1$ converges to both the Gram integral (by tail
squeeze, fully proved) and the Vasyunin formula (by row FTC
decomposition, one convergence axiom), forcing equality. The proof
factors through GCD reduction: $G(j,k) = G(j/d, k/d)$ where
$d = \gcd(j,k)$, so only coprime pairs require the analytic argument.

The per-tile FTC evaluation (\lean{CrossTermFTC.lean}),
the Beatty sequence bound (at most 2 tiles per row for $j \leq k$),
and the Dedekind reciprocity of harmonic tile sums are all
fully proved, establishing that the interaction is
\emph{bounded-range} in the dual lattice.

A further structural result (May~24, 2026) proves that the
\emph{Born approximation} term $E_{\mathrm{ratio}}(j,k) =
\frac{j-k}{2jk}\ln\frac{k}{j}$ is \textbf{always non-positive}
for all $j \neq k$ (\lean{eratio\_entry\_nonpos} in
\lean{GramAssembly.lean}). This follows from the identity
$(a - b) \cdot \ln(b/a) \leq 0$ for all $a, b > 0$
(\lean{sub\_mul\_log\_div\_nonpos}). The Born approximation
never adds energy to the off-diagonal---it always dissipates,
providing structural overcancellation (see~\S\ref{sec:susy}).
\end{observation}

\subsection{Positive Definiteness = Reflection Positivity}

The Cathedral proves that $G$ is positive definite for all $N$
(\lean{gram\_positive\_definite} in \lean{Gram/Bounds.lean}).

\begin{principle}[Reflection Positivity]
In quantum field theory, reflection positivity (Osterwalder--Schrader
axiom RP) ensures that the Hilbert space has positive-definite norm.
The positive definiteness of $G$ is the number-theoretic incarnation
of this axiom: the prime number quantum field satisfies reflection
positivity.

The Cathedral module \lean{White/Kinematics.lean} makes this connection
explicit: it proves the change-of-variables identity
$\int_0^\infty g_N(u)^2\,du = \int_0^1 r_N(x)^2\,dx$
using the diffeomorphism $x = e^{-u}$, which is precisely the
Wick rotation $t \to -iu$ from Minkowski to Euclidean signature.
\end{principle}

